\section{结论与证据资格：哪些演化值得保留}
\label{sec:decision}

\subsection{先给明确判断}
这条研究路线发生了三种性质不同的变化。

\textbf{第一，多解探索转向材料测度与概率体积，是理论深化。} 原先的问题“还有什么材料能解释观测”并未消失。F/F3 已经发现坐标冗余、材料像曲率和非均匀商先验；P3-L/R4 又表明，相容候选、轨迹长度和相对 NLL 改善都不能直接充当后验质量。R5/R6 的原子、重叠原像、条件依赖，正是在修正“应该给谁分配概率”这一基础问题。\refe{H1--H8,R6-A}

\textbf{第二，完整后验转向材料任务，是有实证理由的分叉。} R2、R4 和 R5-B 已独立确认材料平方风险改善；完整后验仍未合格。这个分叉说明“学到有用的条件统计”与“恢复整个概率分布”需要不同资格，不说明完整 posterior 没有价值，也不说明剩余多解不重要。\refe{H6,H7,R6-D}

\textbf{第三，材料任务转向缩短 Native 阶段，目前主要是工程优化。} R6 的费用收益真实存在，但 Direct Probit、强化 Full-q 与计时漂移都削弱了“新的任务最优 Maxwell 原理”这一解释。没有证明概率体积理论导致了这项加速，不能把三个并行结果写成一条已成立的推理链。\refe{R6-C,R6-D}

\begin{tcolorbox}
\textbf{建议：先集中证明一个物理核心，再决定分成两篇。}

优先做非均匀材料上的“表征冗余 / 真正材料变化 / 强方向补偿 / 有限噪声有效阶”区分；随后检验这套几何是否改变可信材料任务积分的成本。当前不建议把 A、B、C 合成一篇总论型 TAP 论文。Paper C 保留为计算模块；Paper B 的风险成像结果是较快的算法基础；Paper A 中的补偿后响应结构具有更强的理论潜力。
\end{tcolorbox}

\subsection{本报告的五种资格标签}
全文把以下情况分开，不把它们统一叫“失败”。

\begin{longtable}{p{.12\linewidth}p{.29\linewidth}p{.51\linewidth}}
\toprule
标签 & 含义 & 典型例子\\
\midrule
T & 命题被证明否定 & 精确边缘独立组合仍可能改变材料；constant-rank 精确核可积，不能称其本质非完整。\\
I & 具体实现未合格 & 当前四段路径未胜直接方向；已有 MALA/SMC/Flow 未达到校准与混合要求。\\
B & 预算、资产或前提未满足 & R6 的四维参考未运行；F2 因冻结驻点门未启动。\\
P & 正面结果尚未取得所需资格 & 局部 GH 加密稳定，但外部模式和尾部无界；有限样本曲线符合幂律。\\
C & 对应冻结合同内已确认 & R2/R4 材料风险收益；R6 原 Full-actual 比较下的两档噪声风险/费用门。\\
\bottomrule
\end{longtable}

T 是否成立取决于精确命题；I 不否定数学对象；B 不是负实验；P 不能当作 C；C 也只覆盖原合同，不能自动延伸到新采集、新网格或更强对照。

\subsection{阅读范围及可追溯性}
R6 指定的主报告、理论、各 Stage、交接、机器汇总与审计均已阅读。另按固定科学版本读回八份 F/L/X 至 R5 历史原报告；A17 的当前最终理论源文件单独核查。R3、R5T 的独立 Pro 全文未在本次附件中取得；对此仅使用用户给出的任务描述和 R5/R6 对旧理论的明确引用，不宣称恢复了缺失全文。

两组 R6 compact NPZ 中全部六臂、共十二个主 MSE，经已暴露结果的零物理复算，与机器汇总逐项一致。此项核查验证摘要转录，不是新独立确认，也不提高参考 posterior 的资格。全部新增 Maxwell RHS 为零。\refe{R6-M}

\section{Research Evolution：研究理想与每次转向的因果关系}
\label{sec:evolution}

\subsection{起点：单幅好图不是全部材料信息}
单幅物理最优解只回答某个目标函数的极小点在哪里。它并未回答：还有多远的材料可以解释数据、这些材料是否属于原先验、它们是否连通、哪类任务在它们之间保持不变，以及不同解释应占多少概率。一个低残差但先验极小的点，和一大片稍高残差但先验质量很大的区域，可以具有完全不同的 posterior 重要性。

最初理想应保留为：\textbf{研究观测仍然允许的不同材料世界，并把几何相容性、概率质量和工程任务的可判断程度分别量化。} 这里“世界”是物理材料，而不是 latent 坐标、current 的任意独立扰动或网络输出像素。

\subsection{L/X/Z：合法材料、局部候选与全局提案}
原始设想是从准确 N3 图附近探索，再用全局提案跨越解族。实际审查表明，N3 在指定原 prior 中未必是合法 anchor：development32 仅两例进 noise-ball，零例满足原相位必要条件。低 MAE 和原 prior 合法性是两回事。任意复像素补偿破坏固定相位时，不能作为原 prior 内多解。\refe{H1}

这否定了具体的候选解释，不是否定多解。径向 proxy 中正确 global-MH 的 event MAE 为 0.004364，优于 local 的 0.012039 和 hybrid 的 0.007243，说明全局提案加正确概率修正有受控正例；复杂 TM 中未找到显著新结构，不能外推为所有问题都没有新结构。

自然产生的问题是：大位移为什么不合法或不相容？是坐标走弯了、先验材料集合弯了，还是 Maxwell 对材料扰动的高阶反馈确实重要？这一分解正是用户原始主线最有价值的部分。

\subsection{F：把坐标曲率、材料像曲率和物理曲率拆开}
沿 latent 曲线 $x(h)$，
\[
\frac{d^2}{dh^2}F(G(x(h)))
=
D^2F[m',m']+DF\,m''.
\]
$D^2G[v,v]$ 又可拆成材料切空间内的坐标加速度和材料像法向的第二基本形式。前者可通过正则重参数化改变；后者反映 prior image 的嵌入几何；$D^2F$ 才是给定材料扰动的物理响应曲率。

F 的 weak/ring 两例中，材料二阶能量约 99.9967\%/85.2285\% 来自坐标项，但映射到数据后各项可能相消，不能把能量比例写成数据响应贡献。旧 Maxwell 二阶向量未保存，只有范数，完整向量归因仍缺失。\refe{H2}

两个受限物理反演得到合法、噪声相容且材料误差很低的点；冻结的驻点阈值未过，F2 未运行。这里有有效正面结果，也有前置门引起的未运行。R7 应纠正研究流程：\textbf{研究局部几何并不总需要 MAP 驻点。} 可先在预先抽定的合法材料上研究响应，随后才检查未知观测反演得到的中心。

\subsection{F3 与 A17-X：商先验已有成果，小 latent 不等于小物理端口}
F3 已消去共同权重尺度 gauge，得到
\[
\rho=\log(w_0/w_1),\qquad
p_\rho(\rho)=
\frac{0.81e^{-|\rho|}-0.16e^{|\rho|}}
{2(0.9-0.4)^2}
\]
在相应合法区间上的非均匀商先验；标签交换和碰合奇点并未消失。故材料商并非 R6 才提出的新对象。\refe{H3}

有限探索获得少数相容端点，但端点间材料距离小于预先定义的任务尺度；直线中点超出 $\Delta\Phi$ 门只否定那些直线路径，不能证明解集不连通。普通 native retraction 在相同费用下不弱于高阶方法，是对该实现的竞争结论，不能永久禁止高阶反馈研究。

A17-X/F3 的审计还显示：七维 latent 生成的材料增量对角算子可有 575 个非零对角元；截成七端口仍损失很大。低参数维度不能推出低更新秩。逐源 dark current 也不自动对应所有源共享的一种材料。\refe{H4}

\subsection{P3-L：近似概率平均可以比物理点估计差得多}
新 P1 使用 Gaussian64 $\times$ Uniform2，quintic 精确平台，24$\times$24 TM/VIE、256 ROI、8 源 $\times$ 32 复接收。128 个独立对象上，特定弱密度 $\times$ strong-Gaussian 近似的材料均值 MSE 为 0.0615021，物理点基线为 0.0110309。11 组 Full reference 未合格，SMC128 退化到一个祖先，strong 条件重要性采样也严重退化。\refe{H5}

这是错误近似导致的风险失败，不是“真正后验均值比点估计差”的反例。平方损失下，真正条件均值仍是 Bayes action。物理代理更准确却未降低 MCSE/时间，也说明目标函数近似误差和概率计算效率不是同一量。由此分出直接材料风险、完整概率和候选探索三条线，有数学必要性。

\subsection{R2：必须同时保留材料风险与多解探索的正面事实}
R2 首批 256 个独立对象中，CNN 将 MSE 从 0.010131196 降至 0.007571318，降幅 25.27\%，调整区间为 22.34\%--27.99\%。第二批确认真实训练谱相对随机谱 decoder 包络的增量约 0.70\%，而反投影独立收益未建立。这意味着主要风险提升不能全部归因于 Maxwell 特有弱空间。\refe{H6}

另一方面，2,048 个合法探索端点中，1,934 个满足 noise-ball 与 $\Delta\Phi\le2$，且有实际 0.02/0.04 材料 RMS 位移。这是值得保留的有限噪声多解证据。弧长平均风险增加约 1.96\%，只说明弧长权重不是已证明正确的 posterior 权重。两个中心的 rank-4 弱 projector 结果给出指定谱截断分布非对合的局部数值证据；它不是全区域的解析证明，也不否定精确 constant-rank fiber。

\subsection{R3 理论到 R4：相对概率学习有效，绝对后验仍不合格}
R3 将问题推向动态几何、非完整控制、Full-66 输运和条件矩。R4 确实执行了相关方法，而非只换命名。

两批新对象中，global statistics + local head 相对 R2 改善材料风险 13.10\%/11.36\%，相对更强 CNN 改善 5.36\%/4.90\%。普通 Gaussian 融合几乎一样好。Full-66 Spline 相对同容量 Gaussian 的 NLL 改善约 5.568/6.171 nat；Maxwell frame 相对随机 frame 改善 1.358/1.602 nat；CFM 的材料 energy score 也改善约 21\%。这些是正面概率学习证据。\refe{H7}

同时，绝对响应校准严重不合格，重要性 ESS 中位数为 1；四段几何路径未获得相同任务距离与费用下的独立优势。两种结论并不冲突。需要特别纠正：
\[
\Phi\sim\Gamma(256,1)
\]
只适用于相应 joint/prior-predictive 跨对象残差校准，不是每个固定 $y$ 条件下 posterior residual 的普遍分布律。

R4 的数值势区间使固定路径重放明显加速，是求解决策的受限正例；它没有改善混合，也不能替代 posterior MCSE。111 个竞争候选也未获得模式质量资格。

\subsection{R5/R5T：响应坐标、混合秩及任务误差}
R5A 表明局部 response-coordinate 积分可改善正则区的材料矩，但常规 Gaussian 在某些案例已足够准确；粗网格会漏掉窄 likelihood。响应半径充分条件不通过，可能是界过松，不能当作逆映射不存在的证明。

R5T/R6 指出 host 原子、不同维数材料像及重叠原像。Full-66 decoder-only host prior 约 0.000688，与截断 S3 的约 0.22086 完全不是同一质量；前者也不是 posterior host mass。\refe{H8,R6-A}

$\sigma^{1/3}$ 与 $\sigma^{4/3}$ 的分离很有启发，但有原子主导、物理首项非零、远分支控制等条件。更关键的是：三次 seam 同时令材料变化与响应变化为 $h^3$；这不等于真正材料线性变化而 Maxwell 只在三阶响应。R7 将在第 \ref{sec:volume} 节严格区分。

\subsection{R5-B 确认与 R6：两条线必须分账}
R5-B 新 512 对象的 Physics / Strong CNN / K32 / Q64 / Distill MSE 分别约为
\[
0.0107810,\ 0.00722154,\ 0.00694331,\ 0.00685473,\ 0.00699702.
\]
风险收益已确认；蒸馏至少保留 80\% 增量仍 HOLD：点估计 80.70\%，调整区间约 [60.28\%,113.12\%]。provider 压缩约 97.97\%，含共同物理 anchor 后整系统推导节时仅 0.0117\%。这些后续确认来自 R6 prior evidence 和审计，不以仍将确认列为未运行的 R5 阶段报告替代。\refe{R6-D}

R6 追求降低物理费用有实际动机，但 A 的概率结构与 C 的优化截断之间没有已证明的因果链。其 6,860,984 RHS、210,355 个状态/LU 与 33 个训练阶段真实存在；不能通过把它们描述成统一框架来提高科学资格。\refe{R6-M}

\subsection{R6 主要结果的并列裁决}
\begin{longtable}{p{.34\linewidth}p{.26\linewidth}p{.32\linewidth}}
\toprule
比较/主张 & 实际结果 & 正确结论\\
\midrule
24+64 q vs Full actual，低/高噪 & MSE 0.00533511 / 0.00955550；原 Full 0.00546356/0.00953433 & 原冻结风险非劣门通过。\\
原窗口完整系统费用 & 节时 34.55\%/32.57\% & 原窗口确认，不是跨窗口保证。\\
同端点 q vs actual 特征 & 低噪风险差 $-1.54137\times10^{-4}$；高噪 $+2.42833\times10^{-5}$ & KKT 独立收益仅低噪确认。\\
Low q vs Full q & 低噪点差 +0.3606\%；1\% NI contrast 上界 $+1.26745\times10^{-5}$ & 强基线低噪 HOLD。高噪仅事后诊断支持。\\
Direct Probit128 & 每噪声 64 个补充对象有竞争力 & 不能与不同分母的全 256 均值直接排名。\\
同窗口重放 & Low 节时 27.65\%/32.77\%；Direct 对 Low 区间均跨零 & 没有 Direct/Low 确定赢家；低噪节时下界 16.82\%。\\
A 的 task 积分 & 局部矩稳定，外部质量未知 & 尚未证明同任务精度下更便宜。\\
\bottomrule
\end{longtable}
上述数据来自冻结预测及相应分析；完整方法族、窗口、置信水平不同，不合并成一个“总体通过”。\refe{R6-D}

\subsection{真正内在统一的地方与尚未统一的地方}
存在内在联系的链是
\[
\text{材料相容几何}
\longrightarrow
\text{该几何上的先验推送与似然重权}
\longrightarrow
\text{材料任务的条件分布及风险}.
\]
计算预算属于“如何近似这些对象”的下一层，只有误差估计器把某次物理误差传递到任务误差时才接上这条链。Native 迭代减少尚未提供这条数学连接。统一研究纲领可以成立；统一论文目前尚不成立。

